\chapter{おわりに}\label{cha:Conclusion}

本研究では、対象とするシステムの特性に合わせた、システム仕様書の章立てを自動生成することを目的として、LLMを用いたシステム仕様書章立て作成のためのプロンプトを提案した。
提案プロンプトは、プロンプトエンジニアリングの技術であるPersona Pattern、One-shot Prompting、Context-Manager Patternを組み合わせることで、標準的な章立ての維持と、対象システムの特性を反映した章立ての拡張の両立を目指して設計した。
提案プロンプトの5つの構成要素を、以下に示す。
\begin{itemize}
    \item 役割付与
    \item ベースの章立て
    \item 開発するシステム情報
    \item システム仕様書の定義
    \item 出力指示
\end{itemize}

提案プロンプトの有用性を確かめるために、BMP画像分割アプリとクラウド型電子契約サービスという2つの異なるシステムを対象として、LLMを用いずに章立てを作成する手法、LLMに自由なプロンプトを与えて章立てを作成する手法、LLMに本研究で提案したプロンプトを与えて章立てを作成する手法の3つの手法で章立てを作成し、その章立ての品質を専門家5人が評価した。
結果として、提案プロンプトは、LLMを用いずに章立てを作成する手法と比較して、章立ての作成に有用であると言えた。
また、提案プロンプトを用いた手法は、LLMに自由なプロンプトを与えた手法と比較して、BMP画像分割アプリのような機能が少ない単純なシステムにおいては、提案プロンプトは有用であると言えた。
一方で、クラウド型電子契約サービスのような機能が多い複雑なシステムにおいては、粒度の適切さや独自性の観点においては、有用であるとは言えなかった。
これは、提案プロンプトに含まれる、章立て生成に対する制約が、複雑なシステムにおいてはLLMの知識活用を阻害し、網羅性が低下したためだと考える。

以下に今後の課題を示す。

\begin{itemize}
    \item より多様なシステムでの提案プロンプトの適用実験と評価\\
    本研究においては、2種類のシステムを対象として実験を行ったが、今後はより多様なシステムを対象として実験を行い、提案プロンプトの有用性と課題を確かめる必要がある。
    \item Chain-of-Thoughtの導入によるシステム仕様書章立て生成プロンプトの改善\\
    本研究の評価実験の結果から、複雑なシステムにおいては、一度に最終的な章立てを生成するのではなく、まず、網羅性を高めることに注力するべきであるという知見を得た。
    そのため、プロンプトにプロンプトエンジニアリング技術の1つであるChain-of-Thoughtを導入し、まずは制約を緩めた状態でシステムに必要な要件や機能を網羅的に列挙させ、その後に、ベースの章立てテンプレートや、システム仕様書の定義に基づいて、章立てを構造化、整形していくというアプローチが有効であると考える。
\end{itemize}

% 本研究では、大規模言語モデル（LLM）を用いて、対象システムの規模や複雑さに関わらず、網羅的かつ論理的なシステム仕様書の章立てを自動生成する手法を提案した。
% 提案手法は、プロンプトエンジニアリングの技術（Persona Pattern, One-shot Prompting, Context-Manager Pattern）を応用し、「役割付与」「ベース章立て」「システム仕様書の定義」「出力指示」といった構成要素を組み合わせることで、標準的な構造の維持と、対象システムの特性に応じた柔軟な拡張の両立を目指して設計された。

% 提案手法の有効性を検証するため、規模の異なる2つのシステム（BMP画像分割アプリ、クラウド型電子契約サービス）を対象に、専門家による比較評価実験を行った。
% 比較対象として、「手動作成」および「自由プロンプトによる生成」を用い、生成された章立ての品質を多角的に評価した。

% 実験の結果、提案手法は手動作成と比較して、「網羅性」「可読性」「構造的整合性」などの多くの項目で高い評価を得た。
% 特に、経験の浅い作成者が陥りやすい章立ての抜け漏れや、階層構造の不整合といった問題を解消し、実用的な品質の章立てを安定して生成できることが示された。
% また、自由プロンプトを用いた手法との比較においては、提案手法の特性として「安定性」と「品質の均一化」が確認された。
% 自由プロンプトは、LLMの知識を最大限に引き出すことで高い網羅性や独自性を発揮する場合がある一方、プロンプト作成者のスキルや偶然性に左右されやすく、冗長な構成や整合性の欠如といったリスクも伴うことが明らかになった。
% これに対し、提案手法は、明示的な制約とテンプレートを与えることで過度な生成を抑制し、評価者から「最も整理されている」「バランスが良い」と評される、実務に適した章立てを安定して出力することに成功した。

% 一方で、課題も明らかになった。
% クラウド型電子契約サービスのような複雑かつ大規模なシステムにおいては、提案手法の制約がLLMの知識活用を一部阻害し、自由プロンプトに比べて網羅性や独自性が低下する傾向が見られた。
% これは、作成者のドメイン知識不足を補うという点においては一定の効果があったものの、大規模システム特有の数多の要件を洗い出し、かつそれを厳密な構造に収めるという高度なタスクにおいては、現在のプロンプト構成では制約が強すぎた可能性がある。

% 今後の課題として、以下の2点が挙げられる。

% 第一に、より多様なシステムに対する適用と評価である。
% 本研究では2種類のシステムのみを対象としたが、さらに多様なドメインや規模のシステムに対して実験を行い、提案手法の汎用性と限界をより詳細に検証する必要がある。

% 第二に、Chain-of-Thought (CoT) の導入によるプロンプトの改善である。
% 複雑なシステムにおいては、一度に最終的な章立てを出力させようとするのではなく、思考プロセスを段階的に踏ませるアプローチが有効であると考えられる。
% 具体的には、まず制約を緩めた状態でシステムに必要な要件や機能を網羅的に列挙させ（発散）、その後にベース章立てや定義に基づいて構造化・整形を行う（収束）という、2段階の生成プロセスを導入することで、網羅性と構造的整合性の両立をより高いレベルで実現できると期待される。


% \section{研究の総括}
% 本研究により得られた成果は以下の通りである。

% 第一に、提案するプロンプトを用いることで、経験の浅い作成者であっても、実用的な品質の章立てを生成できることを確認した。
% 評価実験の結果、提案手法（手法C）は、手動作成（手法A）と比較して、「網羅性」「可読性」「適合性」「構造的整合性」の多くの項目で高い評価を得た。
% 特に、手動作成で見られた「章立ての抜け漏れ」や「階層構造の不整合」といった問題を解消し、安定して高品質な章立てを生成できることが示された。

% 第二に、提案手法は、自由なプロンプト（手法B）と比較して、品質のばらつきを抑制し、バランスの取れた章立てを生成できることを明らかにした。
% 自由なプロンプトを用いた場合、LLMの知識を最大限に引き出すことで、特定のケース（大規模なクラウドサービス等）において高い詳細度や独自性を発揮する可能性がある一方で、冗長な構成や論理的な不整合が生じるリスクも確認された。
% これに対し、提案手法は、明示的な制約とテンプレートを与えることで、過度な生成を抑制し、評価者から「最も整理されている」「バランスが良い」と評される実務に適した章立てを安定して出力することに成功した。

% 以上の結果から、本研究の目的である「標準化とドメイン適応の両立」を実現し、システム仕様書作成の効率化と品質向上に寄与する手法を確立したと言える。